% !TEX program = tectonic
% =============================================================================
%  AULA 3 — Redes sem Fio: Roteamento em Redes Ad Hoc
%  Curso: Redes sem Fio  |  Profa. Anelise Munaretto
%  Protocolos: AODV, DSR, TBRPF, OLSR; Roteamento em redes de sensores.
%  Compilar: tectonic aula3.tex
% =============================================================================
\documentclass[aspectratio=169,11pt]{beamer}
% =============================================================================
%  PREÁMBULO COMÚN para as aulas de Redes sem Fio (Beamer + TikZ)
%  Incluido com % =============================================================================
%  PREÁMBULO COMÚN para as aulas de Redes sem Fio (Beamer + TikZ)
%  Incluido com % =============================================================================
%  PREÁMBULO COMÚN para as aulas de Redes sem Fio (Beamer + TikZ)
%  Incluido com \input{preamble.tex} en cada aula.
%  Paleta de cores e estilos são definidos aqui uma única vez.
% =============================================================================

\usepackage[T1]{fontenc}
\usepackage{amsmath,amssymb,mathtools}
\usepackage{graphicx}
\usepackage{tikz}
\usetikzlibrary{positioning,arrows.meta,calc,backgrounds,decorations.pathmorphing,fit,shadows,shapes.symbols,shapes.geometric}
\usepackage{pgfplots}
\pgfplotsset{compat=1.16}
\usepackage{tabularx,booktabs,multirow,array}
\usepackage[normalem]{ulem}
\usepackage{hyperref}

% ---------- Paleta de cores própria ----------
\definecolor{aneliseAzul}{RGB}{20,60,110}
\definecolor{aneliseVerde}{RGB}{30,140,70}
\definecolor{aneliseLaranja}{RGB}{215,120,20}
\definecolor{aneliseGris}{RGB}{90,90,90}

% ---------- Tema ----------
\usetheme{Madrid}
\usecolortheme{whale}
\setbeamercolor{structure}{fg=aneliseAzul}
\setbeamercolor{title}{fg=aneliseAzul}
\setbeamercolor{frametitle}{fg=aneliseAzul}
\setbeamercolor{block title}{fg=aneliseAzul,bg=aneliseGris!12!white}
\setbeamercolor{block body}{bg=aneliseGris!7!white}
\hypersetup{colorlinks=true,linkcolor=aneliseAzul,urlcolor=aneliseVerde}

% ---------- Comandos auxiliares ----------
\newcommand{\kurose}{Kurose \& Ross, \emph{Computer Networking: a Top-Down Approach}}
\newcommand{\forouzan}{Forouzan, \emph{Data Communications and Networking}}
\newcommand{\stallings}{Stallings, \emph{Wireless Communications \& Networks}}
\newcommand{\tanenbaum}{Tanenbaum, \emph{Computer Networks}}
\newcommand{\rfc}[1]{\href{https://www.rfc-editor.org/rfc/rfc#1.txt}{RFC~#1}}
% -----------------------------------------------------------------------------
%  Estilos TikZ comuns para desenhar redes (posicionamento relativo)
% -----------------------------------------------------------------------------
\tikzset{
  host/.style={draw,rounded corners=2mm,fill=aneliseAzul!15,inner sep=1.5mm,font=\scriptsize},
  rt/.style={draw,rounded corners=1mm,fill=aneliseLaranja!25,inner sep=1mm,font=\scriptsize},
  nube/.style={draw,cloud,aspect=2.2,fill=aneliseGris!15,inner sep=1.5mm,font=\scriptsize},
  el/.style={draw,rounded corners=1.2mm,fill=aneliseGris!18,inner sep=0.8mm,font=\scriptsize},
  enl/.style={thick},
  enlA/.style={thick,-Stealth},
} en cada aula.
%  Paleta de cores e estilos são definidos aqui uma única vez.
% =============================================================================

\usepackage[T1]{fontenc}
\usepackage{amsmath,amssymb,mathtools}
\usepackage{graphicx}
\usepackage{tikz}
\usetikzlibrary{positioning,arrows.meta,calc,backgrounds,decorations.pathmorphing,fit,shadows,shapes.symbols,shapes.geometric}
\usepackage{pgfplots}
\pgfplotsset{compat=1.16}
\usepackage{tabularx,booktabs,multirow,array}
\usepackage[normalem]{ulem}
\usepackage{hyperref}

% ---------- Paleta de cores própria ----------
\definecolor{aneliseAzul}{RGB}{20,60,110}
\definecolor{aneliseVerde}{RGB}{30,140,70}
\definecolor{aneliseLaranja}{RGB}{215,120,20}
\definecolor{aneliseGris}{RGB}{90,90,90}

% ---------- Tema ----------
\usetheme{Madrid}
\usecolortheme{whale}
\setbeamercolor{structure}{fg=aneliseAzul}
\setbeamercolor{title}{fg=aneliseAzul}
\setbeamercolor{frametitle}{fg=aneliseAzul}
\setbeamercolor{block title}{fg=aneliseAzul,bg=aneliseGris!12!white}
\setbeamercolor{block body}{bg=aneliseGris!7!white}
\hypersetup{colorlinks=true,linkcolor=aneliseAzul,urlcolor=aneliseVerde}

% ---------- Comandos auxiliares ----------
\newcommand{\kurose}{Kurose \& Ross, \emph{Computer Networking: a Top-Down Approach}}
\newcommand{\forouzan}{Forouzan, \emph{Data Communications and Networking}}
\newcommand{\stallings}{Stallings, \emph{Wireless Communications \& Networks}}
\newcommand{\tanenbaum}{Tanenbaum, \emph{Computer Networks}}
\newcommand{\rfc}[1]{\href{https://www.rfc-editor.org/rfc/rfc#1.txt}{RFC~#1}}
% -----------------------------------------------------------------------------
%  Estilos TikZ comuns para desenhar redes (posicionamento relativo)
% -----------------------------------------------------------------------------
\tikzset{
  host/.style={draw,rounded corners=2mm,fill=aneliseAzul!15,inner sep=1.5mm,font=\scriptsize},
  rt/.style={draw,rounded corners=1mm,fill=aneliseLaranja!25,inner sep=1mm,font=\scriptsize},
  nube/.style={draw,cloud,aspect=2.2,fill=aneliseGris!15,inner sep=1.5mm,font=\scriptsize},
  el/.style={draw,rounded corners=1.2mm,fill=aneliseGris!18,inner sep=0.8mm,font=\scriptsize},
  enl/.style={thick},
  enlA/.style={thick,-Stealth},
} en cada aula.
%  Paleta de cores e estilos são definidos aqui uma única vez.
% =============================================================================

\usepackage[T1]{fontenc}
\usepackage{amsmath,amssymb,mathtools}
\usepackage{graphicx}
\usepackage{tikz}
\usetikzlibrary{positioning,arrows.meta,calc,backgrounds,decorations.pathmorphing,fit,shadows,shapes.symbols,shapes.geometric}
\usepackage{pgfplots}
\pgfplotsset{compat=1.16}
\usepackage{tabularx,booktabs,multirow,array}
\usepackage[normalem]{ulem}
\usepackage{hyperref}

% ---------- Paleta de cores própria ----------
\definecolor{aneliseAzul}{RGB}{20,60,110}
\definecolor{aneliseVerde}{RGB}{30,140,70}
\definecolor{aneliseLaranja}{RGB}{215,120,20}
\definecolor{aneliseGris}{RGB}{90,90,90}

% ---------- Tema ----------
\usetheme{Madrid}
\usecolortheme{whale}
\setbeamercolor{structure}{fg=aneliseAzul}
\setbeamercolor{title}{fg=aneliseAzul}
\setbeamercolor{frametitle}{fg=aneliseAzul}
\setbeamercolor{block title}{fg=aneliseAzul,bg=aneliseGris!12!white}
\setbeamercolor{block body}{bg=aneliseGris!7!white}
\hypersetup{colorlinks=true,linkcolor=aneliseAzul,urlcolor=aneliseVerde}

% ---------- Comandos auxiliares ----------
\newcommand{\kurose}{Kurose \& Ross, \emph{Computer Networking: a Top-Down Approach}}
\newcommand{\forouzan}{Forouzan, \emph{Data Communications and Networking}}
\newcommand{\stallings}{Stallings, \emph{Wireless Communications \& Networks}}
\newcommand{\tanenbaum}{Tanenbaum, \emph{Computer Networks}}
\newcommand{\rfc}[1]{\href{https://www.rfc-editor.org/rfc/rfc#1.txt}{RFC~#1}}
% -----------------------------------------------------------------------------
%  Estilos TikZ comuns para desenhar redes (posicionamento relativo)
% -----------------------------------------------------------------------------
\tikzset{
  host/.style={draw,rounded corners=2mm,fill=aneliseAzul!15,inner sep=1.5mm,font=\scriptsize},
  rt/.style={draw,rounded corners=1mm,fill=aneliseLaranja!25,inner sep=1mm,font=\scriptsize},
  nube/.style={draw,cloud,aspect=2.2,fill=aneliseGris!15,inner sep=1.5mm,font=\scriptsize},
  el/.style={draw,rounded corners=1.2mm,fill=aneliseGris!18,inner sep=0.8mm,font=\scriptsize},
  enl/.style={thick},
  enlA/.style={thick,-Stealth},
}

\title[Redes sem Fio]{Redes sem Fio\\ Roteamento}
\subtitle{Aula 3}
\author[Anelise Munaretto]{Profa. Anelise Munaretto}
\institute[Curso]{Redes sem Fio --- Ciência da Computação}
\date{2026}

\begin{document}

\begin{frame}[plain]
  \titlepage
  \begin{center}
    \scriptsize Protocolos de roteamento ad hoc: AODV, DSR, TBRPF, OLSR.
  \end{center}
\end{frame}

\section{Características e objetivos}

\begin{frame}{Características das redes sem fio (recursos)}
  \begin{columns}
    \begin{column}{0.5\textwidth}
      \textbf{Recursos limitados}
      \begin{itemize}
        \item quantidade de energia limitada
        \item capacidade de procesamiento limitada (PDAs, sensores,\ldots)
        \item alcance de transmisión pequeno
      \end{itemize}
      \textbf{Topología dinámica}
      \begin{itemize}
        \item mobilidade
        \item falla de nós
        \item nós inactivos em períodos de baixa actividad
      \end{itemize}
    \end{column}
    \begin{column}{0.5\textwidth}
      \textbf{Escalabilidade}
      \begin{itemize}
        \item necessária redundancia --- nós podem falhar
      \end{itemize}
      \textbf{Tempo de vida da rede}
      \begin{itemize}
        \item deve ser o maior possível (principalmente nós sensores)
      \end{itemize}
      \begin{alertblock}{Consecuencia}
        O roteamento deve ser \textbf{eficiente em energia} e \textbf{adaptativo}.
      \end{alertblock}
    \end{column}
  \end{columns}
\end{frame}

\begin{frame}{Objetivos do roteamento ad hoc}
  \begin{block}{Objetivo 1: conectividad total}
    Desarrollar uma infraestructura de rede onde necessário:
    \begin{itemize}
      \item roteamento ad hoc $\rightarrow$ redes auto-organizables, sen mantemento nin administración
      \item roteamento passando por outros terminais
    \end{itemize}
  \end{block}
  \begin{block}{Objetivo 2: mesmo serviço que una rede con infraestructura}
    \begin{itemize}
      \item mobilidade, descubrimiento de servicio, roteamento ad hoc
      \item outros: autoconfiguración, localización, seguridad, consciencia do papel, redes de sensores
    \end{itemize}
  \end{block}
\end{frame}

\begin{frame}{Arquitectura ad hoc}
  \begin{itemize}
    \item creación de uma \textbf{conectividad multisalto} entre um grupo de nós sem fio, frecuentemente en movimento
    \item os nós móviles actúan como \textbf{roteadores} y como \textbf{origen/destino}
    \item necessidade de um protocolo de roteamento que encontre uma ou várias rutas multisalto
      \begin{itemize}
        \item dinámico para adaptarse á topoloxía cambiante
        \item desafíos: interferencias, limitacións de potencia
      \end{itemize}
  \end{itemize}
  \begin{center}
    \begin{tikzpicture}[font=\scriptsize,
        nodo/.style={draw,circle,minimum size=6mm,fill=aneliseAzul!15,inner sep=0.4mm},
        src/.style={draw,circle,minimum size=6mm,fill=aneliseLaranja!35,inner sep=0.4mm},
        dst/.style={draw,circle,minimum size=6mm,fill=aneliseVerde!45,inner sep=0.4mm},
      ]
      \node[src](S){S};
      \node[nodo,right=1.4cm of S](a){a};
      \node[nodo,above=1.2cm of a](b){b};
      \node[nodo,right=1.4cm of a](c){c};
      \node[nodo,right=1.4cm of c](d){d};
      \node[nodo,below=1.2cm of d](e){e};
      \node[dst,right=1.4cm of e](D){D};
      \draw[thick,aneliseVerde] (S) -- node[midway,above=1mm,font=\tiny]{hop 1} (a);
      \draw[thick,aneliseVerde] (a) -- (b);
      \draw[thick,aneliseVerde] (b) -- (c);
      % rota principal: S-a-c-d-e-D
      \draw[thick,aneliseLaranja] (a) -- node[midway,above=1mm,font=\tiny]{hop 2} (c);
      \draw[thick,aneliseVerde] (c) -- (d);
      \draw[thick,aneliseVerde] (d) -- node[midway,above=1mm,font=\tiny]{hop 4} (e);
      \draw[thick,aneliseLaranja] (e) -- (D);
    \end{tikzpicture}
  \end{center}
\end{frame}

\begin{frame}{Necesidades do roteamento ad hoc}
  \begin{columns}
    \begin{column}{0.55\textwidth}
      \textbf{Distribución dos camiños}
      \begin{itemize}
        \item camiños multisalto
        \item sen loops (bucles)
        \item minimizar \emph{overhead}
        \item multicast? inundación eficiente?
      \end{itemize}
    \end{column}
    \begin{column}{0.45\textwidth}
      \begin{itemize}
        \item \textbf{adaptativa á topoloxía dinámica}
        \item \textbf{baixo consumo}: banda e potência
        \item \textbf{escalabilidade} (número de nós)
        \item efectos ante caída de enlaces
      \end{itemize}
    \end{column}
  \end{columns}
\end{frame}

\begin{frame}{Problemas dos protocolos clásicos DV/LS}
  \begin{block}{DV --- Distance Vector (ex.: RIP)}
    \begin{itemize}
      \item podem formar \textbf{bucles (loops)}: moito desperdicio en meios sen fío (banda, potência)
      \item normalmente complexo evitalos
    \end{itemize}
  \end{block}
  \begin{block}{LS --- Link State (ex.: OSPF)}
    \begin{itemize}
      \item custo para almacenar tablas e overhead de comunicación
    \end{itemize}
  \end{block}
  \begin{alertblock}{Aínda máis problemas}
    \begin{itemize}
      \item actualización periódica gasta enerxía (emisor e receptor)
      \item a converxencia debe ser \textbf{rápida en ad hoc}, sen actualizacións moi frecuentes
    \end{itemize}
  \end{alertblock}
\end{frame}

\begin{frame}{Clasificación dos protocolos ad hoc}
  \begin{center}
    \begin{tikzpicture}[font=\scriptsize,
        raiz/.style={draw,rounded corners=1mm,fill=aneliseLaranja!20,inner sep=1mm},
        hoja/.style={draw,rounded corners=1mm,fill=aneliseAzul!15,inner sep=0.8mm,text width=3.2cm,align=center},
        eje/.style={draw,rounded corners=1mm,fill=aneliseGris!18,inner sep=0.8mm},
        flecha/.style={thick,->},
      ]
      \node[raiz](r){Protocolos de roteamento ad hoc};
      \node[eje,below=7mm of r,text width=5cm,align=center](e){Tempo de resposta / consumo};
      \node[hoja,left=2.8cm of e](pro){Proactivos\\(TBRPF, OLSR)};
      \node[hoja,right=2.8cm of e](rea){Reativos\\(AODV, DSR)};
      \draw[flecha] (r) -- (e);
      \draw[flecha] (e.west) to[bend left=10] (pro.north);
      \draw[flecha] (e.east) to[bend right=10] (rea.north);
      \node[draw,fill=aneliseVerde!12,rounded corners=2mm,below=6mm of pro,text width=4cm,align=center](sn){Redes de sensores\\(orientadas a enerxía)};
      \node[draw,fill=aneliseVerde!12,rounded corners=2mm,below=6mm of rea,text width=4cm,align=center](man){MANET\\(orientadas a tempo de resposta)};
      \draw[dashed,->] (pro.south) -- (sn);
      \draw[dashed,->] (rea.south) -- (man);
    \end{tikzpicture}
  \end{center}
\end{frame}

\section{AODV}

\begin{frame}{AODV --- Ad hoc On-Demand Distance Vector}
  \begin{itemize}
    \item Charles Perkins et al. (Nokia Research Center)
    \item \textbf{A protocolo adaptado a redes con alteracións frecuentes de rutas} (manutención de rutas improductiva)
    \item pode adaptarse a redes con gran taxa de transmisión
    \item baseado no \textbf{roteamento de vector de distancia baixo demanda}
      \begin{itemize}
        \item descubrimiento de ruta por \emph{inundación} (flooding)
        \item cálculo de tablas de roteamento (almacenaxe temporal)
        \item avaliación de overheads
      \end{itemize}
  \end{itemize}
  \begin{exampleblock}{RFC}
    RFC 3561 (2003) --- \emph{Ad hoc On-Demand Distance Vector (AODV) Routing}.
  \end{exampleblock}
\end{frame}

\begin{frame}{Mensaxes do protocolo AODV}
  \begin{description}
    \item[Route Request (RREQ)] ``Necesito unha ruta''
    \item[Route Response (RREP)] ``Anunciar a ruta''
    \item[Route Error (RERR)] ``Anular a ruta''
  \end{description}
  \begin{itemize}
    \item respostas periódicas da ruta aos vecinos que actúan como \texttt{HELLO}, para instalación e renovación
  \end{itemize}
  \begin{center}
    \begin{tikzpicture}[font=\scriptsize,
        nodo/.style={draw,circle,minimum size=6mm,fill=aneliseAzul!15,inner sep=0.4mm},
      ]
      \node[nodo](S){A};
      \node[nodo,right=1.4cm of S](f){f};
      \node[nodo,above=1.1cm of f](e){e};
      \node[nodo,right=1.4cm of f](g){g};
      \node[nodo,right=1.4cm of g](h){h};
      \node[nodo,below=1.1cm of h](d){d};
      \node[nodo,right=1.4cm of d](B){B};
      % RREQ flooding
      \draw[->,thick,aneliseLaranja] (S) -- (f);
      \draw[->,thick,aneliseLaranja] (f) -- (g);
      \draw[->,thick,aneliseLaranja] (f) -- (e);
      \draw[->,thick,aneliseLaranja] (g) -- (h);
      \draw[->,thick,aneliseLaranja] (h) -- (d);
      \draw[->,thick,aneliseLaranja] (d) -- (B);
      % RREP de vuelta
      \draw[->,thick,aneliseVerde] (B) to[bend right=20] (d);
      \draw[->,thick,aneliseVerde] (d) to[bend right=20] (h);
      \draw[->,thick,aneliseVerde] (h) to[bend right=20] (g);
      \draw[->,thick,aneliseVerde] (g) to[bend right=20] (f);
      \draw[->,thick,aneliseVerde] (f) to[bend right=20] (S);
      \node[below=1mm of S,font=\tiny]{origen};
      \node[below=1mm of B,font=\tiny]{destino};
    \end{tikzpicture}
  \end{center}
\end{frame}

\begin{frame}{Descubrimiento da ruta (AODV)}
  \begin{columns}
    \begin{column}{0.55\textwidth}
      \begin{enumerate}
        \item Una ruta entre 2 nós é encontrada mediante o envío de \textbf{RREQ} a toda a localidade
        \item inicia en localidade pequena; aumenta con fallas no descubrimiento
        \item un \textbf{RREP} é enviado:
          \begin{itemize}
            \item polo destino
            \item por un nó vecino que xa teña ruta establecida
          \end{itemize}
        \item os nós intermediarios almacenan temporalmente a ruta
      \end{enumerate}
      \textbf{Errores de ruta}
      \begin{itemize}
        \item a ruta expira si non se atualiza
        \item os roteadores almacenan os usuarios da ruta; ao expirar notifícanse (RERR)
        \item nova mensaxe \texttt{route request} será enviada
      \end{itemize}
    \end{column}
    \begin{column}{0.45\textwidth}
      \begin{center}
        \begin{tikzpicture}[font=\scriptsize,
            nodo/.style={draw,circle,minimum size=5mm,fill=aneliseAzul!15,inner sep=0.3mm},
          ]
          \node[nodo](S){A};
          \node[nodo,right=1.3cm of S](b){b};
          \node[nodo,above=1cm of b](c){c};
          \node[nodo,right=1.3cm of b](d){d};
          \node[nodo,right=1.3cm of d](e){e};
          \node[nodo,below=1cm of e](f){f};
          \node[nodo,right=1.3cm of f](B){B};
          \draw[thick,aneliseAzul] (S) -- (b) -- (d) -- (e) -- (B);
          \draw[dashed,thick,aneliseLaranja] (b) -- (c);
          \draw[dashed,thick,aneliseLaranja] (c) -- (d);
          \draw[dashed,thick,aneliseLaranja] (e) -- (f);
          \draw[dashed,thick,aneliseLaranja] (f) -- (B);
          \node[below=0.5mm of S,font=\tiny]{A};
          \node[above=0.5mm of B,font=\tiny]{B};
          \node[draw,fill=aneliseVerde!15,rounded corners=1.5mm,below=6mm of d,text width=2.8cm,align=center](err){ruta A$\rightarrow$B\\expira $\rightarrow$ RERR};
          \draw[dashed,->] (d.south) -- (err);
        \end{tikzpicture}
      \end{center}
    \end{column}
  \end{columns}
\end{frame}

\section{DSR}

\begin{frame}{DSR --- Dynamic Source Routing}
  \begin{itemize}
    \item David Johnson et al. (Rice University)
    \item \textbf{Mensaxes:} Route Discovery, Route Maintenance
    \item \textbf{Múltiples rutas} disponibles
    \item permite a cada emisor \textbf{seleccionar e controlar as rutas utilizadas}
    \item soporta máis de 200 nós
    \item alta mobilidade dos nós
  \end{itemize}
  \begin{exampleblock}{RFC}
    RFC 4728 (2007) --- \emph{The Dynamic Source Routing Protocol for Mobile Ad Hoc Networks (IPv4)}.
  \end{exampleblock}
\end{frame}

\begin{frame}{Descubrimiento da ruta (DSR)}
  \begin{columns}
    \begin{column}{0.5\textwidth}
      \begin{enumerate}
        \item Una ruta entre 2 nós é encontrada mediante \textbf{RREQ}
        \item o RREQ \textbf{constrúe unha ruta cara a fonte} en cada camino (nós rexistran o seu predecesor)
        \item o \textbf{primeiro RREQ que chega é aceptado}; o destino responde por ese camino informando a ruta
        \item o roteamento \textbf{baseado na fonte} é utilizado polo tráfico de datos
      \end{enumerate}
      \begin{exampleblock}{Exemplo}
        Ruta detectada: \texttt{(A, B, F, J, K)} --- encabezada no paquete.
      \end{exampleblock}
    \end{column}
    \begin{column}{0.5\textwidth}
      \begin{center}
        \begin{tikzpicture}[font=\scriptsize,
            nodo/.style={draw,circle,minimum size=5.5mm,fill=aneliseAzul!15,inner sep=0.3mm},
            src/.style={draw,circle,minimum size=5.5mm,fill=aneliseLaranja!35,inner sep=0.3mm},
            dst/.style={draw,circle,minimum size=5.5mm,fill=aneliseVerde!40,inner sep=0.3mm},
          ]
          \node[src](A){A};
          \node[nodo,right=1.3cm of A](B){B};
          \node[nodo,above=1cm of B](C){C};
          \node[nodo,below=1cm of B](D){D};
          \node[nodo,right=1.3cm of B](F){F};
          \node[nodo,right=1.3cm of F](J){J};
          \node[nodo,above=1cm of J](G){G};
          \node[nodo,below=1cm of J](H){H};
          \node[dst,right=1.3cm of J](K){K};
          \node[nodo,above=1.2cm of K](I){I};
          \draw[thick,aneliseLaranja] (A) -- (B) -- (F) -- (J) -- (K);
          \draw[thick,aneliseGris] (B) -- (C);
          \draw[thick,aneliseGris] (B) -- (D);
          \draw[thick,aneliseGris] (J) -- (G);
          \draw[thick,aneliseGris] (J) -- (H);
          \draw[thick,aneliseGris] (K) -- (I);
          \node[below=0.5mm of A,font=\tiny]{fonte};
          \node[above=0.5mm of K,font=\tiny]{destino};
        \end{tikzpicture}
      \end{center}
    \end{column}
  \end{columns}
\end{frame}

\begin{frame}{Confiabilidade do enlace e manutención (DSR)}
  \begin{columns}
    \begin{column}{0.55\textwidth}
      \textbf{Confiabilidade do enlace}
      \begin{itemize}
        \item rutas de tránsito almacenadas (para saber a quem enviar RERR)
        \item confirmacións da capa de rede avaliadas (transferencia, detección de nós con fallas)
        \item tras múltiples retransmisións:
          \begin{itemize}
            \item o nó declarado fóra de servizo
            \item RERR aos emisores que usan ese nó
            \item os paquetes de datos son rexeitados
          \end{itemize}
      \end{itemize}
    \end{column}
    \begin{column}{0.45\textwidth}
      \textbf{Manutención da ruta}
      \begin{itemize}
        \item se un nó roteador ou enlace cambia, as rutas baseadas na fonte deixan de ser válidas
        \item os nós responden con \textbf{RERR}, activando o descubrimiento de nova ruta
        \item os nós poden tentar alterar a ruta e difundila
      \end{itemize}
    \end{column}
  \end{columns}
\end{frame}

\section{TBRPF}

\begin{frame}{TBRPF --- Topology Broadcast Based on Reverse-Path Forwarding}
  \begin{itemize}
    \item Richard Ogier et al. (SRI International)
    \item os nós envían \textbf{só unha parte da tabla de roteamento aos vecinos}
    \item descubrimiento de vecindad con \textbf{HELLO diferenciados}: só se informan cambios de estado dos vecinos
    \item utiliza o algoritmo \textbf{Shortest Path First} (menor camiño) para elección da mellor ruta
  \end{itemize}
  \begin{exampleblock}{RFC}
    RFC 3684 (2004) --- \emph{Topology Dissemination Based on Reverse-Path Forwarding (TBRPF)}.
  \end{exampleblock}
  \begin{block}{Mensaxes}
    \begin{itemize}
      \item \texttt{HELLO}: descubrimiento e manutención da relación entre vecinos
      \item \texttt{TOPOLOGY UPDATE}: distribución dun subconxunto de información para interconectividad
    \end{itemize}
  \end{block}
\end{frame}

\begin{frame}{Distribución da actualización da topoloxía (TBRPF)}
  \begin{itemize}
    \item información da topoloxía é distribuída ao subconxunto de nós que poidan necesitala
      \begin{itemize}
        \item é necesaria información sobre un nó se o roteamento o utiliza\ldots
      \end{itemize}
    \item por iso un nó roteador calcula todas as rutas de todos os nós que posúan tal información
  \end{itemize}
  \begin{center}
    \begin{tikzpicture}[font=\scriptsize,
        nodo/.style={draw,circle,minimum size=5.5mm,fill=aneliseAzul!15,inner sep=0.3mm},
      ]
      \node[nodo](a){a};
      \node[nodo,right=1.3cm of a](b){b};
      \node[nodo,right=1.3cm of b](c){c};
      \node[nodo,right=1.3cm of c](d){d};
      \node[nodo,above=1.2cm of b](e){e};
      \node[nodo,above=1.2cm of c](f){f};
      \draw[thick,aneliseVerde] (a) -- (b) -- (c) -- (d);
      \draw[thick,aneliseGris] (b) -- (e);
      \draw[thick,aneliseGris] (c) -- (f);
      \draw[thick,aneliseGris] (e) -- (f);
      \node[below=0.5mm of b,align=center,text width=2.6cm,font=\tiny]{(subconxunto informado)};
    \end{tikzpicture}
  \end{center}
\end{frame}

\section{OLSR}

\begin{frame}{OLSR --- Optimized Link State Routing}
  \begin{itemize}
    \item Philippe Jacquet et al. (INRIA Rocquencourt)
    \item optimización para redes ad hoc mediante \textbf{MPRs (Multipoint Relays)}
    \item só os nós MPR poden transmitir as mensaxes de descubrimiento da rede
    \item adaptado a redes con \textbf{baixa mobilidade}
    \item utiliza \textbf{Shortest Path First} (menor camiño)
  \end{itemize}
  \begin{exampleblock}{RFC}
    RFC 3626 (2003) --- \emph{Optimized Link State Routing Protocol (OLSR)}.
  \end{exampleblock}
  \begin{block}{Mensaxes}
    \begin{itemize}
      \item \texttt{HELLO}: descubrimiento e manutención de vecinos
      \item \texttt{TC} (Topology Control): distribución de información de conectividade
    \end{itemize}
  \end{block}
\end{frame}

\begin{frame}{Relacións entre nós vecinos e elección de MPR}
  \begin{columns}
    \begin{column}{0.5\textwidth}
      \textbf{Cada equipo emite un \texttt{HELLO} periódico}
      \begin{itemize}
        \item anuncia aos vecinos
        \item determina quen está conectado aínda
        \item selecciona os nós MPR (MultiPoint Relays)
      \end{itemize}
      \textbf{MPRs}
      \begin{itemize}
        \item responsables de transmitir a información da topoloxía
        \item actúan como roteadores entre nós destino
        \item minimizan a retransmisión da información
        \item forman un \textbf{backbone de roteamento}
      \end{itemize}
    \end{column}
    \begin{column}{0.5\textwidth}
      \begin{center}
        \begin{tikzpicture}[font=\scriptsize,
            nodo/.style={draw,circle,minimum size=5mm,fill=aneliseAzul!15,inner sep=0.3mm},
            mpr/.style={draw,circle,minimum size=6mm,fill=aneliseLaranja!35,inner sep=0.3mm},
          ]
          \node[mpr](m1){MPR};
          \node[mpr,right=1.6cm of m1](m2){MPR};
          \node[mpr,right=1.6cm of m2](m3){MPR};
          \node[nodo,above=1.1cm of m1](a){a};
          \node[nodo,above=1.1cm of m3](b){b};
          \node[nodo,below=1.1cm of m2](c){c};
          \draw[thick,aneliseLaranja] (m1) -- (m2) -- (m3);
          \draw[thick,aneliseGris] (m1) -- (a);
          \draw[thick,aneliseGris] (m3) -- (b);
          \draw[thick,aneliseGris] (m2) -- (c);
          \node[below=1mm of m2,font=\tiny]{backbone MPR};
        \end{tikzpicture}
      \end{center}
      \begin{itemize}
        \item a topoloxía pode cambiar: nós móviles, rutas e grupo MPR cambian
      \end{itemize}
    \end{column}
  \end{columns}
\end{frame}

\begin{frame}{Comparación: proactivos vs.\ reativos}
  \begin{columns}
    \begin{column}{0.55\textwidth}
      \textbf{Reativos (baixo demanda) --- AODV, DSR}
      \begin{itemize}
        \item baseados en inundación
        \item non hai garantía de ruta óptima
        \item menor overhead cando hai pouco tráfico
        \item atraso (latencia) no descubrimiento de ruta
      \end{itemize}
    \end{column}
    \begin{column}{0.45\textwidth}
      \textbf{Proactivos --- OLSR, TBRPF}
      \begin{itemize}
        \item descubren continuamente a topoloxía
        \item mensaxes HELLO periódicos
        \item garantía de ruta óptima
        \item posibilidade de obter varias rutas
        \item maior overhead continuo
      \end{itemize}
    \end{column}
  \end{columns}
  \begin{center}
    \begin{tikzpicture}[font=\scriptsize,
        rt/.style={draw,rounded corners=1mm,fill=aneliseAzul!12,inner sep=1mm,text width=2.6cm,align=center},
        re/.style={draw,rounded corners=1mm,fill=aneliseLaranja!15,inner sep=1mm,text width=2.6cm,align=center},
      ]
      \node[rt](p1){Proactivo\\ (OLSR)};
      \node[re,right=3cm of p1](r1){Reativo\\ (AODV)};
      \node[below=3mm of p1,font=\tiny,text width=2.6cm,align=center](pp){ruta sempre dispoñible\\sen atraso de descubrimento};
      \node[below=3mm of r1,font=\tiny,text width=2.6cm,align=center](rr){ruta só cando se pide\\atraso de descubrimento};
    \end{tikzpicture}
  \end{center}
\end{frame}

\section{Roteamento en redes de sensores}

\begin{frame}{Roteamento en redes de sensores: características}
  \begin{columns}
    \begin{column}{0.5\textwidth}
      \textbf{Recursos limitados}
      \begin{itemize}
        \item enerxía limitada (baterías)
        \item capacidade de procesamiento limitada
        \item alcance pequeno
      \end{itemize}
      \textbf{Topoloxía dinámica}
      \begin{itemize}
        \item mobilidade, falla de nós
        \item sensores inactivos en períodos de baixa actividad
      \end{itemize}
    \end{column}
    \begin{column}{0.5\textwidth}
      \textbf{Escalabilidade} e \textbf{tempo de vida}
      \begin{itemize}
        \item redundancia --- nós poden falhar
        \item o obxectivo principal é \textbf{economizar enerxía}
      \end{itemize}
      \textbf{Tipos de roteamento}
      \begin{itemize}
        \item plano (flat), comunicación directa
        \item multisalto de enerxía mínima
        \item aglomeración (clustering), aglomeración dinámica
        \item inundación (flooding), gradiente
        \item roteamento xeográfico
      \end{itemize}
    \end{column}
  \end{columns}
\end{frame}

\begin{frame}{Protocolos de roteamento en sensores}
  \begin{columns}
    \begin{column}{0.5\textwidth}
      \textbf{Plano} (\emph{flat)}
      \begin{itemize}
        \item Difusión Direccionada (Directed Diffusion)
        \item SAR (Sequential Assignment Routing)
        \item SPIN (Sensor Protocols for Information via Negotiation)
        \item Adaptive Local Routing (ALR) / Cooperative Signal Processing
        \item Sensor-MAC, MULTI
      \end{itemize}
    \end{column}
    \begin{column}{0.5\textwidth}
      \textbf{Hierárquico} (\emph{clustering)}
      \begin{itemize}
        \item LEACH
        \item ICA (Inter Cluster Routing)
        \item TEEN
        \item APTEEN (Adaptive TEEN)
        \item SHARP (Hybrid Adaptive Routing)
      \end{itemize}
    \end{column}
  \end{columns}
\end{frame}

\begin{frame}{Difusión Direccionada (Directed Diffusion)}
  \begin{itemize}
    \item proposto por Estrin et al. (1999)
    \item o nó \textbf{nomea os datos} usando un ou máis atributos
    \item o nó cliente \textbf{require os datos} mediante \textbf{intereses}
    \item nós intermediarios \textbf{propagan os intereses}
    \item os intereses establecen \textbf{gradientes de datos} cara ao nó que expresou o interese
  \end{itemize}
  \begin{center}
    \begin{tikzpicture}[font=\scriptsize,
        nodo/.style={draw,circle,minimum size=5.5mm,fill=aneliseAzul!15,inner sep=0.3mm},
        src/.style={draw,circle,minimum size=5.5mm,fill=aneliseLaranja!35,inner sep=0.3mm},
        cli/.style={draw,circle,minimum size=5.5mm,fill=aneliseVerde!45,inner sep=0.3mm},
      ]
      \node[cli](c){cliente};
      \node[nodo,right=1.6cm of c](a){a};
      \node[nodo,right=1.6cm of a](b){b};
      \node[nodo,above=1.2cm of b](d){d};
      \node[src,right=1.6cm of d](f){fonte};
      % interese
      \draw[->,thick,aneliseLaranja] (c) -- node[midway,above=1mm,font=\tiny]{interese} (a);
      \draw[->,thick,aneliseLaranja] (a) -- node[midway,above=1mm,font=\tiny]{} (b);
      \draw[->,thick,aneliseLaranja] (b) -- (d);
      \draw[->,thick,aneliseLaranja] (d) -- (f);
      % gradiente (datos)
      \draw[->,thick,aneliseVerde] (f) to[bend left=15] node[midway,above=1mm,font=\tiny]{datos} (d);
      \draw[->,thick,aneliseVerde] (d) to[bend left=15] (b);
      \draw[->,thick,aneliseVerde] (b) to[bend left=15] (a);
      \draw[->,thick,aneliseVerde] (a) to[bend left=15] (c);
    \end{tikzpicture}
  \end{center}
\end{frame}

\begin{frame}{LEACH --- Low-Energy Adaptive Clustering Hierarchy}
  \begin{itemize}
    \item proposto por Heinzelman et al. (2000)
    \item recolección de datos \textbf{centralizada e periódica} (redes proactivas)
    \item os nós organízanse en \textbf{aglomerados (clusters)}:
      \begin{itemize}
        \item cada cluster ten un \textbf{cluster-head (base local)}
        \item cada nó decide a súa base local (menor custo de comunicación)
        \item \textbf{rotación aleatoria de bases} para nós con maior enerxía
        \item \textbf{agregación de datos local}
      \end{itemize}
    \item cada base crea un \textbf{agendamento TDMA} (rádio apagado, sen colisións)
    \item cada cluster usa un código CDMA diferente
    \item alto consumo nas bases, pero só 5\% son bases; rotación aleatoria $\rightarrow$ \textbf{maior vida da rede}
  \end{itemize}
\end{frame}

\begin{frame}{Comparación: roteamento hierárquico vs.\ plano}
  \begin{center}
    \scriptsize
    \begin{tabular}{@{}p{6.4cm}p{6.4cm}@{}}
      \toprule
      \textbf{Hierárquico} & \textbf{Plano}\\
      \midrule
      evita colisións (TDMA) & non evita colisións\\
      agregación nos cluster-heads & agregación nos nós intermediarios\\
      require sincronización local e global & enlaces formados ``on-the-fly''\\
      overhead de formación de cluster & rutas só nas rexións con datos\\
      disipación uniforme de enerxía & disipación depende do tráfico\\
      alocación xusta de canal & xustiza non garantida\\
      baixa latencia (cluster-heads dispoñibles) & latencia ao esperar nós intermediarios\\
      roteamento simple, non óptimo & roteamento complexo, óptimo\\
      \bottomrule
    \end{tabular}
  \end{center}
\end{frame}

\section{Redes mesh}

\begin{frame}{Roteamento en redes mesh}
  \begin{itemize}
    \item utilizan protocolos de roteamento ad hoc con modificación nas \textbf{métricas}: DSR e OLSR
    \item proxectos de referencia:
      \begin{itemize}
        \item \textbf{MIT Roofnet}: base de roofnet/Meraki mesh networks --- \url{https://pdos.csail.mit.edu/roofnet/}
        \item \textbf{GT-Mesh}: Grupo de Traballo da RNP --- \url{http://mesh.ic.uff.br/}
      \end{itemize}
  \end{itemize}
  \begin{block}{Bibliografía}
    \begin{itemize}
      \item Akyildiz, Wang, ``A survey on wireless mesh networks'', \emph{IEEE Communications Magazine}, Sept.~2005
      \item Akyildiz, Wang, Wang, ``Wireless Mesh Networks: A Survey'', \emph{Computer Networks Journal (Elsevier)}, 2005
    \end{itemize}
  \end{block}
\end{frame}

\section{Bibliografía}

\begin{frame}{Bibliografía}
  \begin{block}{Principal}
    \begin{itemize}
      \item Charles E. Perkins, \emph{Ad Hoc Networking}, Addison-Wesley Professional, 2000
    \end{itemize}
  \end{block}
  \begin{block}{Adicional}
    \begin{itemize}
      \item Royer, Toh, ``A review of current routing protocols for ad hoc mobile wireless networks'', \emph{IEEE Personal Communications}, Abr.~1999
      \item Akyildiz, Wang, Wang, ``Wireless Mesh Networks: A Survey'', \emph{Computer Networks Journal (Elsevier)}, 2005
      \item Liang, Sekercioglu, Mani, ``A survey of multipoint relay based broadcast schemes in wireless ad hoc networks'', \emph{IEEE Communications Surveys \& Tutorials}, 2006
      \item RFCs: AODV (3561), OLSR (3626), TBRPF (3684), DSR (4728)
    \end{itemize}
  \end{block}
\end{frame}

\end{document}